\section{Monte Carlo without Exploring Starts}
    \subsection{Main Idea}
        To deal with the exploring starts, we use $\varepsilon$-greedy function for exploring.
        That is, instead of using the greedy policy for training, we train the $\varepsilon$-greedy policy.
        This little probability of choosing non-greedy action is required, since the greedy action we have depends
        only on current informations. That is, there can be a better choice of action selection, since the action values
        are not perfectly known. Actually, balancing the choice of greedy and non-greedy action is a very
        important problem in reinforcement learning. This problem is called the problem of \textbf{exploitation}
        and \textbf{exploration}.
        \begin{rem}{Problem of Exploitation and Exploration}{ExploitExplor}
            To \textbf{exploit} is to use the current information and select the greedy action.
            To \textbf{explore} is to select the action other than greedy action. Since the true state values cannot be known in
            training, the action that looks greedy may not lead to the highest total reward.
            If the agent exploits too much, the agent may find only suboptimal policies.
            If the agent explores too much, the agent won't be able to find the optimal policy.
            As aforementioned, the balance between exploitation and exploration is critical.
        \end{rem}

    \subsection{Pseudocode}
        \Algorithm{mc-EpsGreedy}
        {Monte Carlo with $\varepsilon$-greedy policy}
        {
            \inputitem{$\pi_0$}{initial policy} \\
            \inputitem{$\varepsilon \in (0,1)$}{probability of exploration}
        }
        {
            \inputitem{$\pi^*$}{optimal policy}
        }
        {
            \tcp{initialize}
            \ForEach{valid state-action pair $(s,a)$}
            {
                $\pi \leftarrow \pi_0$\;
                $q(s,a) \leftarrow$ initial value\;
                \tcp{holds the sum of returns of (sub)episode starting from $(s,a)$}
                $\text{Ret}(s,a) \leftarrow 0$\;
                \tcp{counts the number of visits of $(s,a)$}
                $\text{N}(s,a) \leftarrow 0$\;
            }
            \Repeat{convergence}
            {
                select a starting valid state-action pair $(s,a)$\;
                generate an episode starting from $(s,a)$, following current policy $\pi$: $s_0, a_0, r_1, s_1, a_1, \dots, s_{T-1}, a_{T-1}, r_T$\;
                $G \leftarrow 0$\;
                \For{$t=T-1,T-2,\dots,0$}
                {
                    $G \leftarrow r_{t+1} + \gamma G$\;
                    $\text{Ret}(s,a) \leftarrow \text{Ret}(s,a) + G$\;
                    $\text{N}(s,a) \leftarrow \text{N}(s,a) + 1$\;
                    \tcp{policy evaluation}
                    $q(s_t,a_t) \leftarrow \frac{\text{Ret}(s,a)}{\text{N}(s,a)}$\;
                    \tcp{policy improvement}
                    $\pi_{k+1}(a|s) = \begin{cases}
                        1 - \varepsilon + \frac{\varepsilon}{|\mathcal{A}(s)|} &\text{if } a = \arg\max_{a'}{q(s,a)} \\
                        \frac{\varepsilon}{|\mathcal{A}(s)|} &\text{otherwise}
                    \end{cases}$\;
                }
            }
        }
        
    \subsection{Implementation Details}
        The problem with $\varepsilon$-greedy policy is that it even if the trained policy is optimal in the
        spaces of $\varepsilon$-greedy policies, the policy still has a probability of exploration. To deal with this,
        $\varepsilon$-greedy policy with GLIE (Greedy in the Limit with Infinite Exploration).
        \begin{defn}{GLIE Condition}{glie}
            When following two conditions are satisfied, the policy is called GLIE.
            \begin{itemize}
                \item Infinite Exploration: Every state-action pair must be visited and explored infinitely many times as the training steps go to infinity (i.e., $t \to \infty$).
                \item Limit Greediness: As the learning agent experiences more of the environment, the policy must gradually shift from random exploration to exploiting what it knows, eventually converging entirely to a greedy (optimal) policy.
            \end{itemize}
        \end{defn}
        The most simple way to establish this condition is to use $\varepsilon$-greedy policy with decreasing $\varepsilon$. We can decay the $\varepsilon$ linearly, exponentailly or whatever.
        TODO: search for several GLIE policies.

    \subsection{Pros and Cons}
        Pros:
        \begin{itemize}
            \item Better sample efficiency then MC-Basic;
            \item Does not require exploring starts, and still explores sufficiently;
        \end{itemize}
        Cons:
        \begin{itemize}
            \item Only applicable to episodic tasks, which is a common drawback for all MC methods. This property leads to following problem, too;
            \item High variance since the return holds all the noises and randomnesses, and that return is used for update;
        \end{itemize}